% ============================================================================
%  Rozdział: Struktura modeli sieci neuronowych
%  Plik przeznaczony do samodzielnej kompilacji lub włączenia przez \input{}
%  do głównego dokumentu pracy magisterskiej.
% ============================================================================
\documentclass[12pt,a4paper]{report}
\usepackage[utf8]{inputenc}
\usepackage[T1]{fontenc}
\usepackage[polish]{babel}
\usepackage{amsmath,amssymb}
\usepackage{booktabs}
\usepackage{graphicx}
\usepackage{enumitem}
\usepackage{listings}
\usepackage{xcolor}
\usepackage{hyperref}
\usepackage{geometry}
\geometry{margin=2.5cm}

\lstset{
  basicstyle=\ttfamily\small,
  keywordstyle=\color{blue!70!black},
  commentstyle=\color{green!45!black},
  stringstyle=\color{red!60!black},
  breaklines=true,
  frame=single,
  language=Python,
  showstringspaces=false,
}

\begin{document}

\chapter{Struktura zaimplementowanych modeli sieci neuronowych}
\label{chap:struktura-modeli}

\section{Wprowadzenie}
W katalogu \texttt{src/networks\_tool/models/} zebrano pięć architektur sieci
neuronowych przeznaczonych do regresji parametrów kalibracji kamery. Modele te
dzielą się na dwie rodziny:
\begin{itemize}[leftmargin=*]
  \item \textbf{modele jednoklatkowe} (\emph{single-frame}), które estymują parametry
        kamery na podstawie pojedynczego obrazu wejściowego,
  \item \textbf{modele sekwencyjne} (\emph{sequence}), które agregują informację z krótkiej
        sekwencji klatek przedstawiających tę samą, nieruchomą kamerę obserwującą wzorzec
        z różnych perspektyw.
\end{itemize}

Wszystkie modele realizują wspólny wzorzec projektowy zaczerpnięty z uczenia
transferowego (\emph{transfer learning}): sieć konwolucyjna wstępnie wytrenowana na zbiorze
ImageNet pełni rolę \emph{ekstraktora cech} (ang. \emph{backbone}), z którego usunięto
warstwę klasyfikacyjną, a na jej miejsce dołączono \emph{głowicę regresyjną}
(ang. \emph{regression head}) w postaci wielowarstwowego perceptronu z regularyzacją
metodą \emph{dropout}. Uzasadnieniem tego podejścia jest fakt, że cechy niskiego i
średniego poziomu (krawędzie, łuki, gradienty) wyuczone na naturalnych obrazach są
bezpośrednio powiązane z efektami zniekształceń optycznych, które model ma odtworzyć.

\section{Wspólna definicja wektora wyjściowego}
Modele jednoklatkowe przewidują trójelementowy wektor
\begin{equation}
  \mathbf{y}_{\text{single}} = \big[\, f_{\text{norm}},\; k_1,\; k_2 \,\big],
\end{equation}
gdzie $f_{\text{norm}} = f / W$ jest ogniskową znormalizowaną szerokością obrazu $W$,
a $k_1, k_2$ to współczynniki dystorsji radialnej modelu Browna--Conrady'ego.

Modele sekwencyjne przewidują pełny, dziewięcioelementowy wektor parametrów wewnętrznych
i zniekształceń
\begin{equation}
  \mathbf{y}_{\text{seq}} = \big[\, f_x,\; f_y,\; c_x,\; c_y,\; k_1,\; k_2,\; p_1,\; p_2,\; k_3 \,\big],
\end{equation}
w którym $(f_x, f_y)$ to ogniskowe, $(c_x, c_y)$ współrzędne punktu głównego, a
$(k_1, k_2, k_3)$ oraz $(p_1, p_2)$ odpowiadają współczynnikom dystorsji radialnej i
tangencjalnej. Wartości te są normalizowane rozmiarem obrazu przed treningiem.

\section{Modele jednoklatkowe}

\subsection{ResNet-18 (model bazowy)}
Klasa \texttt{ResNet18CalibrationNet} wykorzystuje sieć ResNet-18 jako ekstraktor cech.
Po globalnym uśrednieniu przestrzennym otrzymywany jest wektor cech o wymiarze
$d = 512$. Głowica regresyjna ma strukturę
\begin{equation}
  512 \xrightarrow{\text{Lin}} 256 \xrightarrow[\text{Dropout }0.3]{\text{ReLU}}
  64 \xrightarrow[\text{Dropout }0.2]{\text{ReLU}} n_{\text{out}} .
\end{equation}
Jest to najlżejszy i najszybszy z wariantów (ok. 11~mln parametrów w części splotowej),
pełniący funkcję punktu odniesienia.

\subsection{ResNet-50}
Klasa \texttt{ResNet50CalibrationNet} zwiększa pojemność modelu, stosując głębszą sieć
ResNet-50, dla której wymiar cech wynosi $d = 2048$. Odpowiednio poszerzono głowicę:
\begin{equation}
  2048 \xrightarrow{\text{Lin}} 512 \xrightarrow[\text{Dropout }0.3]{\text{ReLU}}
  128 \xrightarrow[\text{Dropout }0.2]{\text{ReLU}} n_{\text{out}} .
\end{equation}
Większa liczba parametrów pozwala uchwycić bardziej złożone zależności kosztem
dłuższego czasu wnioskowania.

\subsection{EfficientNet-B0}
Klasa \texttt{EfficientNetB0CalibrationNet} stanowi kompromis pomiędzy jakością a
efektywnością obliczeniową. Ekstraktor EfficientNet-B0 zwraca wektor cech o wymiarze
$d = 1280$, po którym następuje głowica
\begin{equation}
  1280 \xrightarrow{\text{Lin}} 256 \xrightarrow[\text{Dropout }0.3]{\text{ReLU}}
  64 \xrightarrow[\text{Dropout }0.2]{\text{ReLU}} n_{\text{out}} .
\end{equation}

Wspólny schemat przetwarzania modeli jednoklatkowych można zapisać jako
\begin{equation}
  \mathbf{y} = \operatorname{Head}\big(\operatorname{Backbone}(\mathbf{x})\big),
  \qquad \mathbf{x} \in \mathbb{R}^{B \times 3 \times H \times W}.
\end{equation}

\section{Modele sekwencyjne}

\subsection{CNN + LSTM}
Klasa \texttt{CNNLSTMCalibrationNet} przetwarza tensor o kształcie
$(B, T, 3, H, W)$, gdzie $T$ oznacza liczbę klatek w oknie czasowym. Każda klatka jest
niezależnie kodowana wspólnym ekstraktorem ResNet-18, a otrzymany wektor cech jest
rzutowany warstwą liniową do przestrzeni ukrytej o wymiarze $256$:
\begin{equation}
  \mathbf{h}_t = \operatorname{ReLU}\!\big(W_{\text{proj}}\,\operatorname{Backbone}(\mathbf{x}_t)\big),
  \qquad t = 1, \dots, T .
\end{equation}
Sekwencja wektorów cech $\{\mathbf{h}_t\}$ jest następnie przetwarzana przez
dwuwarstwową sieć LSTM (rozmiar stanu ukrytego $256$, \emph{dropout} $0.2$):
\begin{equation}
  (\mathbf{o}_1, \dots, \mathbf{o}_T) = \operatorname{LSTM}(\mathbf{h}_1, \dots, \mathbf{h}_T).
\end{equation}
Do warstwy regresyjnej trafia stan z ostatniego kroku czasowego $\mathbf{o}_T$:
\begin{equation}
  \mathbf{y} = \operatorname{Head}(\mathbf{o}_T),
  \qquad \operatorname{Head}:\; 256 \to 128 \to n_{\text{out}} .
\end{equation}
Model ten jest przeznaczony do sytuacji, w których parametry zniekształceń są stałe w
obrębie sekwencji, a zmienia się jedynie perspektywa obserwacji wzorca.

\subsection{CNN + Transformer}
Klasa \texttt{CNNTransformerCalibrationNet} zastępuje rekurencyjny agregator mechanizmem
samouwagi (\emph{self-attention}). Cechy każdej klatki są rzutowane liniowo do przestrzeni
osadzeń o wymiarze $d_{\text{model}} = 256$, a następnie wzbogacane o uczone kodowanie
pozycyjne $\mathbf{p}_t$:
\begin{equation}
  \mathbf{z}_t = W_{\text{proj}}\,\operatorname{Backbone}(\mathbf{x}_t) + \mathbf{p}_t .
\end{equation}
Sekwencja $\{\mathbf{z}_t\}$ jest przetwarzana przez enkoder transformera złożony z
$3$ warstw, każda z $8$ głowicami uwagi i warstwą \emph{feed-forward} o wymiarze
$4 d_{\text{model}} = 1024$. Reprezentacja sekwencji powstaje przez uśrednienie po
wymiarze czasowym (\emph{mean pooling}):
\begin{equation}
  \bar{\mathbf{z}} = \frac{1}{T} \sum_{t=1}^{T} \operatorname{Encoder}(\mathbf{z})_t,
  \qquad \mathbf{y} = \operatorname{Head}(\bar{\mathbf{z}}).
\end{equation}
Mechanizm uwagi pozwala modelować bezpośrednie relacje pomiędzy dowolnymi klatkami, co
czyni ten wariant odpowiednim dla dłuższych sekwencji.

\section{Funkcja straty}
Wszystkie modele trenowane są z użyciem funkcji straty \emph{log-cosh}
(klasa \texttt{LogCoshLoss}), stanowiącej gładkie przybliżenie straty $L_1$:
\begin{equation}
  \mathcal{L}(\mathbf{y}, \hat{\mathbf{y}}) =
  \frac{1}{n} \sum_{i=1}^{n} \log\!\big(\cosh(y_i - \hat{y}_i)\big).
\end{equation}
Dla małych błędów zachowuje się jak strata kwadratowa $L_2$, a dla dużych --- jak
$L_1$, co zapewnia odporność na wartości odstające. Rozwiązanie to zaczerpnięto z pracy
DeepCalib (Bogdan i in., 2018).

\section{Zestawienie architektur}
Tabela~\ref{tab:modele} podsumowuje kluczowe własności wszystkich pięciu wariantów.

\begin{table}[htbp]
  \centering
  \caption{Zestawienie architektur modeli z katalogu \texttt{models/}.}
  \label{tab:modele}
  \small
  \begin{tabular}{@{}llccl@{}}
    \toprule
    \textbf{Model} & \textbf{Backbone} & \textbf{Wymiar cech} & \textbf{Wejście} & \textbf{Agregacja} \\
    \midrule
    ResNet18CalibrationNet      & ResNet-18       & 512  & $(B,3,H,W)$   & --- \\
    ResNet50CalibrationNet      & ResNet-50       & 2048 & $(B,3,H,W)$   & --- \\
    EfficientNetB0CalibrationNet& EfficientNet-B0 & 1280 & $(B,3,H,W)$   & --- \\
    CNNLSTMCalibrationNet       & ResNet-18       & 512  & $(B,T,3,H,W)$ & LSTM (2 warstwy) \\
    CNNTransformerCalibrationNet& ResNet-18       & 512  & $(B,T,3,H,W)$ & Transformer (3 warstwy) \\
    \bottomrule
  \end{tabular}
\end{table}

\section{Sposób użycia}
Modele importuje się bezpośrednio z pakietu, a wyboru wariantu dokonuje się na podstawie
pola \texttt{model\_name} w konfiguracji treningu:

\begin{lstlisting}
from src.networks_tool.models import (
    ResNet18CalibrationNet,
    ResNet50CalibrationNet,
    EfficientNetB0CalibrationNet,
    CNNLSTMCalibrationNet,
    CNNTransformerCalibrationNet,
)
\end{lstlisting}

Modele jednoklatkowe przyjmują tensor o kształcie $(B, 3, H, W)$, natomiast modele
sekwencyjne --- tensor $(B, T, 3, H, W)$. Domyślny rozmiar obrazu wejściowego wynosi
$224 \times 224$ pikseli, zgodnie z wymaganiami wstępnie wytrenowanych ekstraktorów cech.

\end{document}
